% Section 5.5: 本章小结

本章在第四章方法口径下，按「设置 $\to$ 结果 $\to$ 分析」三段依次推进，以 H1（Base + Direct，pass@1 = 52.40\%）为共同基线，得到以下实证结论：
\begin{itemize}
    \item \textbf{训练侧主效应（pass@1）}：$H4-H1 = +8.1$\,pp、$H7-H1 = +4.8$\,pp、$H9-H1 = +6.1$\,pp。纯代码 LoRA 增益最大且在 pass@1/5/10 三档同号正向，CoT 监督 LoRA 在 Direct 推理下回落（$H7-H4 = -3.3$\,pp）说明 CoT 监督对 Direct 推理构成格式开销；
    \item \textbf{推理侧主效应（pass@1）}：$H2-H1 = +1.7$\,pp、$H3-H1 = +5.9$\,pp、$H3-H2 = +4.2$\,pp。0-shot CoT 推理在基座上仅微弱正向，few-shot CoT 推理则带来显著边际增益；
    \item \textbf{训练 $\times$ 推理交互}：$(H8-H7)-(H2-H1) = +5.0$\,pp、$(H10-H9)-(H3-H1) = +3.1$\,pp 显著为正，$(H5-H4)-(H2-H1) = -0.1$\,pp、$(H6-H4)-(H3-H1) = -0.6$\,pp 在 0 附近——\textbf{训练侧是否显式监督 CoT 格式是触发「LoRA 放大 CoT」的必要条件}：监督 CoT 时形成正协同，仅监督代码时退化为正交可加；
    \item \textbf{few-shot 一致性}：$H3-H2 = +4.2$\,pp 与 $H6-H5 = +3.7$\,pp 符号一致且量级差 $\le 0.6$\,pp，判定 few-shot 为\textbf{训练侧无关的结构性增益}；$H10-H8 = +3.6$\,pp 与前两条量级相当，未显著放大；
    \item \textbf{CoT 等价性}：$H7-H8 = -6.7$\,pp、$H9-H10 = -9.0$\,pp，两组不等式方向均与「训练侧等价替代」假设相反——\textbf{「CoT 进参数」与「CoT 进推理」是互补关系}；
\end{itemize}

\textbf{主要结论}：本工作在 KodCode 评测集 1\,000 题、$T=0.3$、$n=10$ 采样的口径下，证实了 LoRA 在算法代码生成任务上具备稳定的训练侧主效应（pass@1 最高提升 $+8.1$\,pp 达 60.50\%），并通过 10 组析因实验厘清了 LoRA 与 CoT 的交互边界——\textbf{两者并非简单可加，而是受训练侧是否监督 CoT 格式调节，且训练侧 CoT 监督需与推理侧 CoT 两阶段同分布触发才能取得最大协同（H10 = 67.50\%，相对 H1 提升 $+15.1$\,pp）}。该结论对应 Ch1.4 贡献 (3)，并构成第六章 6.1 节「主要结论」与 6.3 节「未来工作」的实证依据。